% Options for packages loaded elsewhere
% Options for packages loaded elsewhere
\PassOptionsToPackage{unicode}{hyperref}
\PassOptionsToPackage{hyphens}{url}
\PassOptionsToPackage{dvipsnames,svgnames,x11names}{xcolor}
%
\documentclass[
  letterpaper,
  DIV=11,
  numbers=noendperiod]{scrartcl}
\usepackage{xcolor}
\usepackage{amsmath,amssymb}
\setcounter{secnumdepth}{-\maxdimen} % remove section numbering
\usepackage{iftex}
\ifPDFTeX
  \usepackage[T1]{fontenc}
  \usepackage[utf8]{inputenc}
  \usepackage{textcomp} % provide euro and other symbols
\else % if luatex or xetex
  \usepackage{unicode-math} % this also loads fontspec
  \defaultfontfeatures{Scale=MatchLowercase}
  \defaultfontfeatures[\rmfamily]{Ligatures=TeX,Scale=1}
\fi
\usepackage{lmodern}
\ifPDFTeX\else
  % xetex/luatex font selection
\fi
% Use upquote if available, for straight quotes in verbatim environments
\IfFileExists{upquote.sty}{\usepackage{upquote}}{}
\IfFileExists{microtype.sty}{% use microtype if available
  \usepackage[]{microtype}
  \UseMicrotypeSet[protrusion]{basicmath} % disable protrusion for tt fonts
}{}
\makeatletter
\@ifundefined{KOMAClassName}{% if non-KOMA class
  \IfFileExists{parskip.sty}{%
    \usepackage{parskip}
  }{% else
    \setlength{\parindent}{0pt}
    \setlength{\parskip}{6pt plus 2pt minus 1pt}}
}{% if KOMA class
  \KOMAoptions{parskip=half}}
\makeatother
% Make \paragraph and \subparagraph free-standing
\makeatletter
\ifx\paragraph\undefined\else
  \let\oldparagraph\paragraph
  \renewcommand{\paragraph}{
    \@ifstar
      \xxxParagraphStar
      \xxxParagraphNoStar
  }
  \newcommand{\xxxParagraphStar}[1]{\oldparagraph*{#1}\mbox{}}
  \newcommand{\xxxParagraphNoStar}[1]{\oldparagraph{#1}\mbox{}}
\fi
\ifx\subparagraph\undefined\else
  \let\oldsubparagraph\subparagraph
  \renewcommand{\subparagraph}{
    \@ifstar
      \xxxSubParagraphStar
      \xxxSubParagraphNoStar
  }
  \newcommand{\xxxSubParagraphStar}[1]{\oldsubparagraph*{#1}\mbox{}}
  \newcommand{\xxxSubParagraphNoStar}[1]{\oldsubparagraph{#1}\mbox{}}
\fi
\makeatother


\usepackage{longtable,booktabs,array}
\newcounter{none} % for unnumbered tables
\usepackage{calc} % for calculating minipage widths
% Correct order of tables after \paragraph or \subparagraph
\usepackage{etoolbox}
\makeatletter
\patchcmd\longtable{\par}{\if@noskipsec\mbox{}\fi\par}{}{}
\makeatother
% Allow footnotes in longtable head/foot
\IfFileExists{footnotehyper.sty}{\usepackage{footnotehyper}}{\usepackage{footnote}}
\makesavenoteenv{longtable}
\usepackage{graphicx}
\makeatletter
\newsavebox\pandoc@box
\newcommand*\pandocbounded[1]{% scales image to fit in text height/width
  \sbox\pandoc@box{#1}%
  \Gscale@div\@tempa{\textheight}{\dimexpr\ht\pandoc@box+\dp\pandoc@box\relax}%
  \Gscale@div\@tempb{\linewidth}{\wd\pandoc@box}%
  \ifdim\@tempb\p@<\@tempa\p@\let\@tempa\@tempb\fi% select the smaller of both
  \ifdim\@tempa\p@<\p@\scalebox{\@tempa}{\usebox\pandoc@box}%
  \else\usebox{\pandoc@box}%
  \fi%
}
% Set default figure placement to htbp
\def\fps@figure{htbp}
\makeatother





\setlength{\emergencystretch}{3em} % prevent overfull lines

\providecommand{\tightlist}{%
  \setlength{\itemsep}{0pt}\setlength{\parskip}{0pt}}



 


% Page furniture for the project PDFs, ported from the exam-documentclass originals the 2024 bootcamp used before the Quarto migration.
%
% Pulled in by every project's YAML:
%
%   format:
%     pdf:
%       include-in-header: header.tex
%     latex:
%       include-in-header: header.tex
%
% Both keys are needed: pdf and latex are separate formats to Quarto, and with the key on pdf alone the shipped .tex comes out without the header its own .pdf has.
%
% The exam class provided \header, \headrule and \footer; Quarto's pdf output uses the default article class, so fancyhdr supplies the same furniture.
% The exam class's \numpages becomes \pageref*{LastPage}, which needs two passes --- Quarto runs latexmk, so that is already handled.
% The star matters: plain \pageref is a hyperref link, so the page count comes out link-blue.
%
% Editing this file changes every project PDF, so it is folded into deploy.sh's GLOBAL_HASH: a change here rebuilds the whole site rather than serving stale PDFs.
%
% ---------------------------------------------------------------------------
% CURRENTLY INERT: every rule below is commented out, matching the P230A repo
% this was ported from. Uncomment the block to turn the running header on.
% ADAPTING THIS FOR ANOTHER COURSE: change the two strings below, nothing else.
% ---------------------------------------------------------------------------

% \usepackage{fancyhdr}
% \usepackage{lastpage}

% \pagestyle{fancy}
% \fancyhf{}

% \fancyhead[L]{MCSB Bootcamp: Mathematical and Computational Track}  % <- course name
% \fancyfoot[L]{\textit{jun.allard@uci.edu}}                          % <- contact

% \fancyfoot[R]{Page \thepage\ of \pageref*{LastPage}}

% \renewcommand{\headrulewidth}{0.4pt}
% \renewcommand{\footrulewidth}{0pt}

% % Quarto's title block starts the document with \maketitle, which forces the plain pagestyle on the first page and would otherwise drop the header exactly where it is most wanted.
% \fancypagestyle{plain}{%
%   \fancyhf{}%
%   \fancyhead[L]{MCSB Bootcamp: Mathematical and Computational Track}%
%   \fancyfoot[L]{\textit{jun.allard@uci.edu}}%
%   \fancyfoot[R]{Page \thepage\ of \pageref*{LastPage}}%
%   \renewcommand{\headrulewidth}{0.4pt}%
%   \renewcommand{\footrulewidth}{0pt}%
% }
\KOMAoption{captions}{tableheading}
\makeatletter
\@ifpackageloaded{caption}{}{\usepackage{caption}}
\AtBeginDocument{%
\ifdefined\contentsname
  \renewcommand*\contentsname{Table of contents}
\else
  \newcommand\contentsname{Table of contents}
\fi
\ifdefined\listfigurename
  \renewcommand*\listfigurename{List of Figures}
\else
  \newcommand\listfigurename{List of Figures}
\fi
\ifdefined\listtablename
  \renewcommand*\listtablename{List of Tables}
\else
  \newcommand\listtablename{List of Tables}
\fi
\ifdefined\figurename
  \renewcommand*\figurename{Figure}
\else
  \newcommand\figurename{Figure}
\fi
\ifdefined\tablename
  \renewcommand*\tablename{Table}
\else
  \newcommand\tablename{Table}
\fi
}
\@ifpackageloaded{float}{}{\usepackage{float}}
\floatstyle{ruled}
\@ifundefined{c@chapter}{\newfloat{codelisting}{h}{lop}}{\newfloat{codelisting}{h}{lop}[chapter]}
\floatname{codelisting}{Listing}
\newcommand*\listoflistings{\listof{codelisting}{List of Listings}}
\makeatother
\makeatletter
\makeatother
\makeatletter
\@ifpackageloaded{caption}{}{\usepackage{caption}}
\@ifpackageloaded{subcaption}{}{\usepackage{subcaption}}
\makeatother
\usepackage{bookmark}
\IfFileExists{xurl.sty}{\usepackage{xurl}}{} % add URL line breaks if available
\urlstyle{same}
\hypersetup{
  pdftitle={Project 4: Phosphorylation--dephosphorylation},
  pdfauthor={Jun Allard},
  colorlinks=true,
  linkcolor={blue},
  filecolor={Maroon},
  citecolor={Blue},
  urlcolor={Blue},
  pdfcreator={LaTeX via pandoc}}

\ifXeTeX
\special{pdf:trailerid [ <12ff9ff09db876af9dda6aef1448b9ba> <12ff9ff09db876af9dda6aef1448b9ba> ]}
\fi
\ifPDFTeX
\pdftrailerid{}
\pdftrailer{/ID [ <12ff9ff09db876af9dda6aef1448b9ba> <12ff9ff09db876af9dda6aef1448b9ba> ]}
\fi
\ifLuaTeX
\pdfvariable trailerid {[ <12ff9ff09db876af9dda6aef1448b9ba> <12ff9ff09db876af9dda6aef1448b9ba> ]}
\fi


\title{Project 4: Phosphorylation--dephosphorylation}
\usepackage{etoolbox}
\makeatletter
\providecommand{\subtitle}[1]{% add subtitle to \maketitle
  \apptocmd{\@title}{\par {\large #1 \par}}{}{}
}
\makeatother
\subtitle{MCSB Bootcamp --- Mathematical and Computational Track}
\author{Jun Allard}
\date{}
\begin{document}
\maketitle


\subsection{Background and Part I}\label{background-and-part-i}

Many proteins are activated (phosphorylated) by enzymes called kinases,
and deactivated (dephosphorylated) by enzymes called phosphatases.
Intuitively, the fraction of active protein is determined by the balance
of kinases to phosphatases. This balance is maintained by constant
activation and inactivation in what is called a \emph{futile cycle}.

\begin{enumerate}
\def\labelenumi{\alph{enumi}.}
\tightlist
\item
  Sketch a diagram where the horizontal axis is the ratio of kinase to
  phosphatase (say, in log scale), and the vertical axis is the fraction
  of activated protein in steady state. Use your intuition to sketch
  what you think the curve should look like.
\end{enumerate}

In mathematics, we call this plot a bifurcation diagram (the steady
state as a function of a parameter). In biology, we call this particular
plot a dose-response curve (how the output of a system responds to the
dose of an input).

\subsection{Model statement}\label{model-statement}

To activate the protein, the kinase must first bind to the inactive
form, then catalyze activation. It may also unbind before it has a
chance to catalyze. The same is true for the phosphatase. There are
therefore 6 species of molecule. Let \([A]\) be the concentration of A,
\([AP]\) be the concentration of the AP complex, and so on. We have
defined six populations, with concentrations:

\[
\begin{aligned}
{}[A] &  & &\text{unbound active protein}\\
[I] &  & &\text{unbound inactive protein}\\
[AP] &  & &\text{complex of active protein and phosphatase}\\
[P] &= [P_\mathrm{tot}] - [AP] & &\text{unbound phosphatase}\\
[IK] &  & &\text{complex of inactive protein and kinase}\\
[K] &= [K_\mathrm{tot}] - [IK] & &\text{unbound kinase}
\end{aligned}
\]

Assume the total concentration of phosphatase \([P_\mathrm{tot}]\) and
kinase \([K_\mathrm{tot}]\) are constant. Therefore, we do not need to
track the concentration of P, since it is determined by
\([P_\mathrm{tot}]\) and \([AP]\). Similarly for K. Therefore, there are
4 concentrations to keep track of.

There are 6 chemical reactions, shown below.

\[
\mathrm{I} + \mathrm{K}
\;\underset{k^I_\mathrm{off}}{\overset{k^I_\mathrm{on}}{\rightleftharpoons}}\;
\mathrm{IK}
\;\xrightarrow{\;k^A_\mathrm{cat}\;}\;
\mathrm{A} + \mathrm{K}
\]

\[
\mathrm{A} + \mathrm{P}
\;\underset{k^A_\mathrm{off}}{\overset{k^A_\mathrm{on}}{\rightleftharpoons}}\;
\mathrm{AP}
\;\xrightarrow{\;k^I_\mathrm{cat}\;}\;
\mathrm{I} + \mathrm{P}
\]

Let's assume reactions are proportional to the concentration of the
reactants, so \(\mathrm{A}+\mathrm{P} \rightarrow \mathrm{AP}\) occurs
at rate
\(k^A_\mathrm{on} [P] [A] = k^A_\mathrm{on} ([P_\mathrm{tot}] - [AP]) [A]\).
Therefore, we have defined 6 rate parameters:

\[
\begin{aligned}
k^A_\mathrm{on} &= 10 \ \mathrm{s}^{-1}\mu\mathrm{M}^{-1} & \qquad k^I_\mathrm{on} &= 10 \ \mathrm{s}^{-1}\mu\mathrm{M}^{-1}\\
k^A_\mathrm{off} &= 10 \ \mathrm{s}^{-1} & \qquad k^I_\mathrm{off} &= 10 \ \mathrm{s}^{-1}\\
k^A_\mathrm{cat} &= 100 \ \mathrm{s}^{-1} & \qquad k^I_\mathrm{cat} &= 10 \ \mathrm{s}^{-1}
\end{aligned}
\]

\subsection{Part II}\label{part-ii}

\begin{enumerate}
\def\labelenumi{\alph{enumi}.}
\setcounter{enumi}{1}
\item
  Write down 4 ODEs describing the populations of A, I, AP and IK. The
  first two equations are

  \[
   \begin{aligned}
   \frac{d[A]}{dt} &= -k^A_\mathrm{on} \left( [P_\mathrm{tot}] - [AP]\right) [A] + k^A_\mathrm{off} [AP] + k^A_\mathrm{cat} [IK],\\
   \frac{d[AP]}{dt} &= +k^A_\mathrm{on} \left( [P_\mathrm{tot}] - [AP]\right) [A] - k^A_\mathrm{off} [AP] - k^I_\mathrm{cat} [AP].
   \end{aligned}
   \]
\item
  Assume that at time \(t=0\), all of the protein is in the inactive,
  unbound state, and that its concentration is
  \(I_\mathrm{tot}=1~\mu\mathrm{M}\). What are the appropriate initial
  conditions for the ODEs?
\item
  Write code to simulate this ODE system to steady state. Use the
  parameters stated above, the initial conditions from the previous
  question, and \(I_\mathrm{tot}=1~\mu\mathrm{M}\),
  \(K_\mathrm{tot}=1~\mu\mathrm{M}\) and
  \(P_\mathrm{tot}=1~\mu\mathrm{M}\).
\end{enumerate}

\subsection{Part III}\label{part-iii}

\begin{enumerate}
\def\labelenumi{\alph{enumi}.}
\setcounter{enumi}{4}
\tightlist
\item
  Do a parameter sweep in \(K_\mathrm{tot}\), logarithmic from
  \(K_\mathrm{tot} = 10^{-3}~\mu\mathrm{M}\) to
  \(K_\mathrm{tot} = 10^{+2}~\mu\mathrm{M}\). Plot the total amount of
  activated protein in steady state as a function of \(K_\mathrm{tot}\)
  (with a log axis in \(K_\mathrm{tot}\)). Does the plot match the
  intuitive sketch you made in question (a)?
\item
  Now do the same parameter sweep in \(K_\mathrm{tot}\), identical to
  the previous question except \(I_\mathrm{tot}=100~\mu\mathrm{M}\).
  What changed? In what circumstances might a cell gain an advantage
  from this kind of dose-response curve (compared to the previous
  question with \(I_\mathrm{tot}=1~\mu\mathrm{M}\))? In what
  circumstances might a cell be disadvantaged by a dose-response curve
  like this (compared to the previous question with
  \(I_\mathrm{tot}=1~\mu\mathrm{M}\))?\footnote{This type of behavior in
    systems biology is called \emph{ultrasensitivity}. This particular
    type of ultrasensitivity is called \emph{Goldbeter-Koshland
    ultrasensitivity}.}
\end{enumerate}




\end{document}
